% 对应题号 M1-10；图与本题条件、推导、结论逐项对应。
% 图契约：印刷宽16cm；最小字号8.4pt；几何/场或曲线在左，计算关系在右。
\documentclass[tikz,border=7pt]{standalone}
\usepackage{fontspec,xeCJK,amsmath,amssymb,graphicx}
\IfFontExistsTF{Times New Roman}{\setmainfont{Times New Roman}}{\setmainfont{TeX Gyre Termes}}
\IfFontExistsTF{Heiti SC}{\setCJKmainfont[AutoFakeBold=1.5]{Heiti SC}}{\setCJKmainfont[AutoFakeBold=1.5]{FandolHei-Regular}}
\usetikzlibrary{arrows.meta,bending,calc,patterns,positioning}
\definecolor{ink}{HTML}{243745}
\definecolor{blue}{HTML}{275DAB}
\definecolor{green}{HTML}{177C63}
\definecolor{red}{HTML}{BD4935}
\definecolor{orange}{HTML}{AA6C16}
\definecolor{purple}{HTML}{715393}
\tikzset{every node/.style={font=\fontsize{9.3}{12.5}\selectfont,text=ink},
 axis/.style={-{Stealth[length=4pt]},line width=.65pt,draw=ink!65},
 vec/.style={-{Stealth[length=5pt]},line width=1pt,draw=blue},
 short/.style={-{Stealth[length=3pt]},line width=.8pt},
 guide/.style={dashed,draw=ink!40,line width=.5pt},
 card/.style={rounded corners=2pt,draw=blue!65,fill=blue!4,align=center,text width=6.6cm,inner xsep=9pt,inner ysep=7pt},
 title/.style={font=\fontsize{10.5}{14}\selectfont\bfseries,text width=15cm,align=center},
 label/.style={font=\fontsize{8.4}{11}\selectfont,align=center}}
\begin{document}
\begin{tikzpicture}[x=1cm,y=1cm]
\path[use as bounding box] (-8,-3.4) rectangle(8,3.4);
\node[title] at (0,3) {M1-10\quad 变化的标量因子产生附加散度};
\begin{scope}[shift={(-4,0)}]
\draw[axis] (-2.4,-1.1)--(2.5,-1.1) node[right] {$x$};
\draw[axis] (0,-2.25)--(0,1.1) node[above] {$y$};

\draw[blue,line width=1pt] (-.8,-1.9)--(2.3,1.2) node[above] {$f=x$};
\draw[vec,green] (.4,.5)--(.7,.5);\draw[vec,green] (1.2,.5)--(2,.5);
\node[label] at (0,2.2) {$\mathbf A=\mathbf e_x$；箭长示意 $x>0$ 的乘积场};
\node[label] at (0,-2.75) {$\mathbf A$ 均匀，但 $f\mathbf A=x\mathbf e_x$ 随位置变化};

\end{scope}\node[card,draw=blue!65,fill=blue!4] at (3.95,1.65) {$\nabla\cdot\mathbf A=0$，$\nabla f=\mathbf e_x$};
\node[card,draw=green!65,fill=green!4] at (3.95,0) {$\nabla\cdot(f\mathbf A)=\partial_x x=1$};
\node[card,draw=orange!65,fill=orange!4] at (3.95,-1.65) {$f\nabla\cdot\mathbf A+(\nabla f)\cdot\mathbf A$\\$=0+1=1$};
\end{tikzpicture}\end{document}
